\textbf{Descripción y correspondencia de vistas arquitectónicas.}
La descripción de arquitectura aplica ISO/IEC/IEEE 42010 al sistema compuesto por canales, tres núcleos, cinco procesos especializados, Runtime Operacional de Borde y plataformas de datos. Los interesados son Gerencia y Administración, Torre y Marinería, Escuela, Varadero, contratistas, armadores y apoderados, y los responsables de Arquitectura, Seguridad y Operación de Synaptix. Sus preocupaciones se agrupan en continuidad y custodia, integridad y privacidad, tiempos operacionales, recuperación, mantenibilidad y gobierno. Cada vista utiliza los mismos nombres e identificadores de SD3 y SD5 \parencites{iso42010}[RT-02.03, p.~7]{pucv-transversal}.

La Tabla~\ref{tab:vistas-42010} individualiza los cinco puntos de vista, su preocupación y su representación. Las figuras existentes son modelos de la oferta; el esquema de datos de la Figura~\ref{fig:datos-arquitectura} completa su correspondencia.

\begin{table}[htbp]
\centering
\fontsize{9}{11}\selectfont
\begin{tabularx}{\linewidth}{|l|X|X|}
\hline
\EncabezadoTabla{Vista} & \EncabezadoTabla{Preocupación} & \EncabezadoTabla{Modelo y ubicación}\\ \hline
VL: lógica & Límites, capacidades y contratos & Figuras~\ref{fig:4-2}, \ref{fig:4-3} y \ref{fig:4-4}\\ \hline
VP: procesos & Captura, concurrencia y sincronización & Figura~\ref{fig:4-5}; contratos y deduplicación de 4.1\\ \hline
VDP: despliegue & Localización, falla y recuperación & Figuras~\ref{fig:4-7}, \ref{fig:4-9}, \ref{fig:4-10} y \ref{fig:4-14}\\ \hline
VD: datos & Propiedad, persistencia y evidencia & Figura~\ref{fig:datos-arquitectura}; SD5, 5.1 y 5.2\\ \hline
VS: seguridad & Identidad, fronteras y privilegios & Figura~\ref{fig:4-6}; controles de 4.1 y 4.2\\ \hline
\end{tabularx}
\caption{Puntos de vista y modelos de arquitectura}
\label{tab:vistas-42010}
\FuenteTabla{Descripción de Synaptix conforme a ISO/IEC/IEEE 42010 y BT RT-02.03.}
\end{table}

VL identifica componentes, dominios y dependencias permitidas; VP utiliza secuencias de captura, entrega y confirmación para mostrar quién ejecuta y cuándo persiste; VDP ubica sus procesos y almacenes en regiones, zonas y recinto; VD representa autoridades y flujos de datos, y VS atraviesa cada frontera con identidad, autorización, protección y evidencia. Ninguna vista sustituye a las demás. La restricción de no consultar esquemas ajenos se comprueba en VL y VD; la confirmación durable antes del ACK se comprueba en VP y VD; la continuidad necesita coherencia entre VP, VDP y VS. Toda diferencia conserva registro de problema, vistas afectadas, decisión y solución antes de liberar una nueva línea base.

La Figura~\ref{fig:datos-arquitectura} sitúa los registros transaccionales, las proyecciones y la evidencia, y muestra las autoridades que impiden que una copia de lectura confirme una operación.

\begin{figure}[htbp]
\centering
\setstretch{1}
\begin{tikzpicture}[
  bloque/.style={draw=SynNavyLight,fill=SynSurface,align=left,
    font=\fontsize{9.5}{11.5}\selectfont,text width=64mm,inner sep=2mm,
    minimum height=28mm},
  banda/.style={bloque,text width=136mm,minimum height=0mm},
  flujo/.style={-{Latex},draw=SynNavyLight,thick},
  copia/.style={flujo,dashed}]
  \node[banda,fill=SynNavy,text=SynWhite] (propiedad)
    {\textbf{Propiedad por nueve dominios}\\
    Cada dominio escribe su esquema; API y eventos intercambian referencias.
    Seguridad registra identidad, consentimiento y decisión.};
  \node[bloque,below=7mm of propiedad.south west,anchor=north west] (local)
    {\textbf{PostgreSQL local + SQLite en terreno}\\
    Hechos críticos, RLH, secuencias y outbox.
    Autoridad local de puesto y captura durable.
    Terminal aislado no arbitra una amarra.};
  \node[bloque,below=7mm of propiedad.south east,anchor=north east] (cloud)
    {\textbf{PostgreSQL transaccional cloud}\\
    Esquemas de los núcleos y Seguridad.
    Maestros comerciales y hechos incorporados.
    Inbox y restricciones de identidad.};
  \node[banda,below=7mm of local.south west,anchor=north west] (intercambio)
    {\textbf{Sincronización y eventos versionados}\\
    Hechos local a cloud; maestros firmados cloud a local.
    Idempotencia, secuencia, validación y persistencia antes del ACK.
    El evento transporta información y conserva su propietario.};
  \node[bloque,below=7mm of intercambio.south west,anchor=north west] (lectura)
    {\textbf{PostgreSQL de proyecciones}\\
    Recursos y permisos separados.
    Vistas reconstruibles, versión y antigüedad.
    Sin escritura sobre maestros ni confirmación crítica.};
  \node[bloque,below=7mm of intercambio.south east,anchor=north east] (evidencia)
    {\textbf{S3: evidencia y Parquet}\\
    Original, versiones, huella y retención.
    Referencias desde el hecho propietario.
    Athena consulta historia autorizada.};
  \node[banda,below=7mm of lectura.south west,anchor=north west] (proteccion)
    {\textbf{Protección y recuperación}\\
    Cifrado, permisos por dominio y campo sensible; réplica regional y copias.
    La restauración conserva identidades, relaciones y trazabilidad.
    Una réplica o proyección no sustituye el respaldo ni la autoridad.};
  \draw[flujo] (propiedad.south -| local.north) -- (local.north);
  \draw[flujo] (propiedad.south -| cloud.north) -- (cloud.north);
  \draw[flujo] (local.south) -- (intercambio.north -| local.south);
  \draw[flujo] (cloud.south) -- (intercambio.north -| cloud.south);
  \draw[copia] (intercambio.south -| lectura.north) -- (lectura.north);
  \draw[flujo] (intercambio.south -| evidencia.north) -- (evidencia.north);
  \draw[copia] (lectura.south) -- (proteccion.north -| lectura.south);
  \draw[copia] (evidencia.south) -- (proteccion.north -| evidencia.south);
\end{tikzpicture}
\caption{Vista arquitectónica de datos, autoridades y proyecciones}
\label{fig:datos-arquitectura}
\FuenteFigura{Elaboración de Synaptix desde SD4, 4.1, y SD5, 5.1/5.2;
  \parencites{iso42010}[RT-02.03, p.~7]{pucv-transversal}.
  Las flechas continuas representan contratos; las discontinuas, derivación o copia.}
\end{figure}


Los hechos críticos nacen en su autoridad local y se incorporan al propietario cloud por entrega durable, mientras los maestros comerciales viajan en sentido inverso como versiones. La evidencia documental se relaciona por identificador y huella, sin almacenar binarios dentro de cada transacción. Las proyecciones reciben eventos y conservan su versión; una consulta reconstruida no adquiere autoridad. El detalle de entidades, atributos, clasificación y retención se mantiene en SD5 y sus anexos, con correspondencia explícita a estos almacenes.

El registro contractual de decisiones se entrega en el archivo \ReferenciaAnexo{Anexos de SD4}{SYNAPTIX-Subdocumento4-Anexos.pdf}. Sus registros DA-01 a DA-20 individualizan fecha, contexto, elección, alternativas, fundamento y consecuencias. Arquitectura conserva el registro versionado en el repositorio del CLIENTE durante implementación y operación. Una modificación añade una decisión sucesora y relación de sustitución; no reescribe silenciosamente el antecedente. Seguridad y Operación revisan su impacto y la contraparte autorizada decide los cambios de alcance, mediante control contractual. Los ADR se actualizan en cada cambio y se entregan con la línea base y el informe mensual; su fecha de formulación no representa una aprobación del CLIENTE \parencite[RT-02.04, p.~7]{pucv-transversal}.
